\chapter*{贯穿案例：从复位向量到一台可扩展的操作系统}
\addcontentsline{toc}{chapter}{贯穿案例：从复位向量到一台可扩展的操作系统}

本册使用一套真实的 x86-64 教学系统作为贯穿案例。案例不借助现成固件或第三方
引导器把内核直接放进内存，而是从处理器复位向量接管机器：自研 ROM 初始化最早
的串口和磁盘访问，装入 Stage 1；Stage 1 建立长模式，校验并装入 ELF64 内核；
内核再接管描述符表、异常、内存、中断和设备。这样，每个操作系统对象都能回到
一个已经发生的机器状态变化，而不是停在术语定义上。

\begin{keyidea}
项目主线负责回答“当前机器真实拥有什么、下一步为什么失败、最小实现怎样通过
证据验收”；其后的原理正文负责把同一问题扩展到多任务、多核、并发、故障恢复和
现代系统规模。案例给出骨架，深入章节解释骨架继续生长时必须遵守的规则。
\end{keyidea}

\section*{先确定案例的边界}

模拟器在案例中只承担硬件角色：它提供 x86-64 处理器、内存、串口、IDE、
计时器、中断控制器和键盘等器件。复位后执行什么字节、怎样读取磁盘、何时打开
分页、如何解析 ELF、异常入口保存什么、哪个中断需要确认，均由目标系统自行
决定。这个边界很重要：如果外部环境已经替内核装入 ELF、建立页表或准备
BootInfo，读者就看不到这些对象为什么存在，也无法区分硬件事实和软件策略。

案例当前已经完成 v0.0 至 v1.0 的十三阶段路线，共有 \(99\) 个进入目标系统
的 C++ 或汇编源文件，约 \(14208\) 行非空、非纯注释目标代码。代码量不是完成
证据；更重要的事实是，机器已经能够从 \code{0xFFFFFFF0} 取到自研指令，最终
建立四个独立进程、真实 PIT 抢占、条件等待、管道、持久文件系统和普通 Ring 3
Shell。Shell 的输入来自真实 PS/2 IRQ1，文件能够跨两个全新 QEMU 实例保留，
损坏的 superblock 会在进入用户态以前被拒绝。

v1.0 已经形成 README、roadmap、四份连续 ADR 与 v0.9--v1.0 release 共同确认
的完成边界。v0.9 把单次用户执行装入 PCB、独立 PML4、Ring 0 栈和
176 字节保存帧；v0.10 增加 Blocked、具名等待和 64 字节管道；v0.11 建立
固定磁盘格式、写回缓存与 Dirty/Clean 提交；v1.0 再用八槽描述符表统一控制台、
管道、文件和目录，并使无 Ready、仍有 Blocked 的机器进入可由中断唤醒的
\code{sti; hlt; cli} 空闲路径。

最近一次完整验证通过 \(73/73\) 项测试。七份用户 ELF 分别接受 AMD64、入口、
PT\_LOAD、W\^X 和运行时依赖审计；单元、集成与固定种子随机测试覆盖调度、
管道、控制台、描述符、解析器和文件系统不变量。正常 QEMU 路径逐字输入
\(109\) 个字符，要求 submitted 与 read 相等、dropped 与 buffered 均为零，
同时保留四进程退出、256 字节管道守恒、文件同步和地址空间隔离。持久化测试在
同一可写磁盘上连续启动两个全新实例，再破坏 superblock 并验证第三次启动不会
格式化、进入 Ring 3 或输出 \code{READY}。因此 v1.0 是由构造、采用、失败和
跨层证据共同闭合的教学系统基线。

取得内核控制权以前，五个里程碑依次建立可验证的启动链：

\begin{longtable}{|L{0.12\linewidth}|L{0.23\linewidth}|L{0.31\linewidth}|L{0.24\linewidth}|}
\toprule
版本 & 当前机器新增能力 & 必须保留的状态或契约 & 成功证据 \\
\midrule
v0.0 & 固定目标平台与空镜像基线 & CPU 型号、内存大小、设备集合与有界运行时间 & 空镜像不能伪造启动成功 \\\tablerowrule
v0.1 & ROM 从复位向量执行并输出串口 & 复位地址、16 位入口、栈、COM1 寄存器与轮询预算 & \code{RESET}、\code{SERIAL_READY} \\\tablerowrule
v0.2 & 固件读取并交接 Stage 1 & 磁盘描述符、LBA、加载区间、入口和负载校验 & \code{STAGE1_HEADER_VALID}、\code{STAGE1_LOADED} \\\tablerowrule
v0.3 & Stage 1 进入 64 位长模式 & A20、GDT、页表、CR0、CR3、CR4 与 EFER 的提交顺序 & \code{LONG_MODE} \\\tablerowrule
v0.4 & ELF64 内核获得控制权 & Kernel 容器、PT\_LOAD、BSS、BootInfo 与调用 ABI & \code{KERNEL_TRANSFER}、\code{KERNEL ENTERED} \\
\bottomrule
\end{longtable}

内核取得控制权以后，后九个里程碑把临时启动环境逐步替换成内核长期拥有的
控制结构，并最终形成可交互、可持久化的用户环境：

\begin{longtable}{|L{0.12\linewidth}|L{0.23\linewidth}|L{0.31\linewidth}|L{0.24\linewidth}|}
\toprule
版本 & 当前机器新增能力 & 必须保留的状态或契约 & 成功证据 \\
\midrule
v0.5 & 内核接管异常控制边界 & GDT、TSS、IDT、统一异常帧、IST 与 panic 现场 & \code{BREAKPOINT_HANDLED} \\\tablerowrule
v0.6 & 内核拥有内存与页表 & 物理内存图、页帧状态、四级页表、W\^X、guard page 与早期堆 & \code{MEMORY_PERMISSIONS_VALID} \\\tablerowrule
v0.7 & 内核接收异步设备事件 & PIC 路由、IRQ 入口、EOI、PIT tick、键盘状态与 ATA 事务 & \code{TIMER_SELF_TEST_PASSED} 与键盘事件 \\\tablerowrule
v0.8 & 内核执行受硬件隔离的用户程序 & 用户 ELF、U/S 页、TSS.RSP0、\code{INT 0x80} 与用户故障收束 & \code{HELLO_FROM_RING3}、退出与隔离测试 \\\tablerowrule
v0.9 & 四进程可被真实 PIT 抢占 & PCB、独立 PML4、每进程内核栈、CR3/TSS.RSP0 切换与页回收 & \code{ADDRESS_SPACE_ISOLATED}、抢占统计与资源回收 \\\tablerowrule
v0.10 & 进程能够按条件阻塞和通信 & Blocked、等待原因、acquire/release 锁、64 B 管道与端点关闭 & 256 B 守恒、EOF 与 block/wakeup 统计 \\\tablerowrule
v0.11 & 文件能够跨启动保存并拒绝损坏 & 固定磁盘格式、位图、inode、目录、LRU 缓存与 Dirty/Clean & 同盘双启动、全盘一致性与损坏停止 \\\tablerowrule
v1.0 & 普通用户程序获得交互环境 & 八槽描述符、控制台 FIFO、目录迭代、空闲唤醒与 Shell & 十条命令、109 字符零丢失、73 项回归 \\
\bottomrule
\end{longtable}

\section*{一条不能跳步的控制权链}

复位到内核不是一次函数调用。每一阶段都只能依赖上一阶段明确交付的状态，并且
必须在交出控制权以后停止修改已转移的对象。下图中的箭头因此表示所有权和 ABI
边界，而不只是“接下来执行哪个函数”。

\begin{pdfdiagram}
\node[os control, text width=2.2cm] (reset)
  {CPU 复位\\CS 隐藏基址 + IP\\物理地址 \code{0xFFFFFFF0}};
\node[os state, text width=2.2cm, right=7mm of reset] (rom)
  {128 KiB ROM\\16 位入口、串口\\ATA PIO};
\node[os compute, text width=2.2cm, right=7mm of rom] (stage)
  {Stage 1\\A20、GDT、页表\\长模式};
\node[os memory, text width=2.2cm, right=7mm of stage] (loader)
  {ELF64 装载\\PT\_LOAD、BSS\\BootInfo};
\node[os state, text width=2.2cm, right=7mm of loader] (kernel)
  {freestanding 内核\\异常、内存\\中断与设备};
\draw[os strong bus] (reset) -- node[os label, above] {取指} (rom);
\draw[os strong bus] (rom) -- node[os label, above] {远转移} (stage);
\draw[os strong bus] (stage) -- node[os label, above] {两遍装载} (loader);
\draw[os strong bus] (loader) -- node[os label, above] {RDI + CALL} (kernel);
\end{pdfdiagram}
\diagramnote{案例的主控制权链。每个阶段都要定义入口处理器状态、可使用的内存、
失败后的停机语义和下一阶段入口；缺少任一项，后续成功只能算偶然现象。}

这条链把“装入”“接受”“运行”和“内核就绪”分成不同事件。ROM 读完 Stage 1
扇区，只能说明字节已经进入约定内存；校验通过后执行远转移，才说明 Stage 1
接受控制。Stage 1 把 ELF 段复制到目标地址，只能说明映像已经建立；BootInfo
准备完成并按 ABI 调用入口，才说明内核获得控制。内核输出 \code{READY} 之前，
还必须自行验证交接状态并建立长期拥有的结构。

\section*{v0.1：处理器怎样找到第一条项目指令}

案例 ROM 大小为 128 KiB，映射在物理地址
\code{0xFFFE0000..0xFFFFFFFF}。复位时，处理器第一次取指位于
\code{0xFFFFFFF0}；这个地址对应 ROM 文件中的 \code{0x1FFF0}。该位置保存
一条 16 位 near jump，把 IP 从 \code{0xFFF0} 改为 \code{0xF000}，但不改变
CS 的隐藏基址，因此目标物理地址是 \code{0xFFFFF000}。

\begin{longtable}{|L{0.21\linewidth}|L{0.23\linewidth}|L{0.25\linewidth}|L{0.21\linewidth}|}
\toprule
坐标 & 数值 & 谁解释它 & 它不能单独证明什么 \\
\midrule
ROM 文件偏移 & \code{0x1FFF0} & 镜像布局与地址译码 & CPU 此刻是否已经取指 \\\tablerowrule
线性/物理地址 & \code{0xFFFFFFF0} & 复位架构状态与平台映射 & 跳转目标是否有合法代码 \\\tablerowrule
入口偏移 & \code{0xF000} & 16 位 near jump 与当前 CS 基址 & 栈和数据段是否可用 \\\tablerowrule
低端栈顶 & \code{0x7000} & ROM 入口软件 & 该区域是否会与后续装载重叠 \\
\bottomrule
\end{longtable}

入口首先关闭可屏蔽中断并清除方向标志，再建立 DS、ES、SS、SP 和 BP。这个
顺序不能由 C++ 运行时替代，因为此时没有谁承诺栈、BSS、全局构造或 ABI 已经
成立。串口初始化也采用有界 THRE 轮询；设备若永久不就绪，固件进入稳定停机，
而不是在未知状态下继续读盘。

\begin{lstlisting}
firmware_entry():
    // 复位没有提供可直接信任的软件栈，先建立最小 16 位执行环境。
    disable_interrupts()
    clear_direction_flag()
    set_data_segments(0)
    set_stack(segment=0, offset=0x7000)

    // 串口是后续所有阶段共用的最早证据通道；等待必须有上限。
    if not initialize_com1_with_bounded_poll():
        halt_forever()
    write_marker("RESET")
    write_marker("SERIAL_READY")
\end{lstlisting}

\section*{v0.2：从“读到扇区”到“接受一个可执行阶段”}

固件不能把 LBA 0 的任意 512 字节直接当成程序。案例先读取一份固定宽度、
自描述的 Stage 1 描述符，逐字段校验 magic、版本、头长度、加载段、入口偏移、
扇区数、标志、负载 LBA 和校验值。描述符通过后，固件才读取负载扇区到
\code{0x00008000} 起始的窗口。

\begin{pdfdiagram}
\node[os memory, text width=3.1cm] (disk0)
  {LBA 0\\Stage 1 描述符\\格式、范围、校验};
\node[os memory, text width=3.1cm, below=12mm of disk0] (diskn)
  {LBA 1..N\\Stage 1 负载\\按扇区读取};
\node[os compute, text width=3.2cm, right=24mm of disk0] (verify)
  {ROM 验证器\\先校验描述符\\再校验完整负载};
\node[os state, text width=3.2cm, right=24mm of diskn] (ram)
  {RAM \code{0x8000..}\\负载字节与入口\\校验前不可执行};
\draw[os strong bus] (disk0) -- node[os label, above] {ATA PIO} (verify);
\draw[os strong bus] (diskn) -- node[os label, below] {256 个 16 位 word/扇区} (ram);
\draw[os strong bus] (verify) -- node[os label, right] {范围成立后允许装入} (ram);
\end{pdfdiagram}
\diagramnote{Stage 1 的格式事实与负载字节分开校验。读盘完成不等于内容可执行，
只有描述符、加载范围和完整负载共同通过验证，固件才交出 CS:IP。}

ATA PIO 状态也不能压成一个 \code{ready} 布尔量。每轮要先判断 BSY 是否清零，
再区分 ERR、DF 与 DRQ；轮询预算耗尽和设备报告错误具有不同故障原因。成功后，
DATA 端口每次返回 16 位，精确读取 256 次才得到一个 512 字节扇区。短读、
多读或忽略错误位都会破坏描述符边界。

\begin{lstlisting}
read_one_sector(lba, destination):
    // 命令字段必须在通知设备以前全部写好。
    select_primary_master_lba28(lba)
    write_sector_count(1)
    write_lba_fields(lba)
    issue_read_sectors()

    for attempt in 0..POLL_LIMIT:
        status = read_status()
        // ERR/DF 是稳定失败；BSY 只表示当前不能取数。
        if status has (ERR or DF):
            return DEVICE_ERROR
        if status has DRQ and not status has BSY:
            // DATA 端口宽度为 16 位，一个扇区恰好读取 256 次。
            read_256_words(destination)
            return SUCCEEDED
    return TIMEOUT
\end{lstlisting}

校验失败不会回到调用者继续执行未知代码。案例为串口失败、ATA 永久忙、设备
错误、描述符损坏和负载损坏分别保留失败终态。失败终态的价值在于：测试不仅
能观察“没有成功”，还可以判断系统在哪一条契约上拒绝了输入。

\section*{v0.3：长模式是按依赖顺序提交的处理器状态}

Stage 1 进入时仍处于 16 位实模式。它必须先打开并验证 A20，再装载平坦 GDT，
设置 \code{CR0.PE} 并通过远跳转刷新代码段，才能在 32 位保护模式建立进入
IA-32e 所需的页表和控制状态。案例使用三页结构建立 PML4、PDPT 和 PD，并以
32 个 2 MiB 大页对低 64 MiB 做恒等映射。

\begin{pdfdiagram}
\node[os state, text width=2.25cm] (real) {16 位实模式\\A20 已验证};
\node[os compute, text width=2.25cm, right=9mm of real] (protected)
  {GDT + \code{CR0.PE}\\远跳到 32 位};
\node[os memory, text width=2.25cm, right=9mm of protected] (tables)
  {PML4/PDPT/PD\\低 64 MiB 恒等映射};
\node[os compute, text width=2.25cm, right=9mm of tables] (enable)
  {\code{CR4.PAE}\\\code{CR3}\\\code{EFER.LME}\\\code{CR0.PG}};
\node[os state, text width=2.25cm, right=9mm of enable] (long)
  {64 位子模式\\\code{EFER.LMA}=1};
\draw[os strong bus] (real) -- (protected);
\draw[os strong bus] (protected) -- (tables);
\draw[os strong bus] (tables) -- (enable);
\draw[os strong bus] (enable) -- node[os label, above] {L 位代码段远跳} (long);
\end{pdfdiagram}
\diagramnote{长模式转换的依赖顺序。任一控制位或页表只在前置状态已经成立时才
具有预期含义；每一步回读验证后再推进，失败则停在当前阶段。}

“已经写过寄存器”不是证据。案例在每一步回读相应位，并在分页开启后检查
\code{IA32_EFER.LMA}，证明处理器实际进入了长模式。日志
\code{PAE_READY}、\code{LME_READY}、\code{PAGING_ENABLED} 和
\code{LONG_MODE} 因而对应不同提交边界，不能用最后一条日志反推所有中间状态
都必然正确。

\section*{v0.4：ELF 装载是校验事务，不是边读边跳}

Kernel 容器与 Stage 1 占用不同磁盘区域。LBA 65 保存 Kernel 描述符，
LBA 66 起保存精确长度的 \code{kernel.elf}，文件末尾到扇区边界的填充必须为
零。Stage 1 先验证描述符和完整负载 CRC32，再解析 ELF64。

解析采用两遍算法。第一遍检查 ELF 头和全部程序头，确认 PT\_LOAD 的文件范围、
目标范围、4 KiB 对齐、权限、W\^X、段间不重叠与入口可执行；这一遍不改目标
内存。只有所有段共同成立，第二遍才复制文件字节并清零
\code{p_memsz - p_filesz} 对应的 BSS。于是任一段失败都不会留下“前几个段已经
覆盖内存、后一个段才报错”的半提交映像。

\begin{longtable}{|L{0.18\linewidth}|L{0.28\linewidth}|L{0.28\linewidth}|L{0.16\linewidth}|}
\toprule
阶段 & 读取的状态 & 允许修改的状态 & 失败结果 \\
\midrule
容器验证 & 描述符、文件长度、扇区数、CRC32 与零填充 & 暂存区与验证状态 & 目标段保持不变 \\\tablerowrule
ELF 第一遍 & ELF 头、程序头、段范围、权限和入口 & 仅验证账本 & 拒绝整个映像 \\\tablerowrule
ELF 第二遍 & 已通过验证的 PT\_LOAD 集合 & 目标段和 BSS & 进入交接准备 \\\tablerowrule
BootInfo 交接 & 内存布局、内核栈、文件长度与段数 & 固定宽度交接对象 & 不调用内核 \\
\bottomrule
\end{longtable}

BootInfo v2 为 104 字节，字段均为固定宽度的 64 位值。它至少表达 magic、
版本、结构长度、内核文件长度、加载段数、初始页表根、恒等映射上限、内核栈顶
以及物理内存图的地址、数量和条目宽度。Stage 1 把 BootInfo 地址放入 RDI，
在满足 System V AMD64 入口对齐的栈上调用内核入口。

\begin{pdfdiagram}
\node[os memory, text width=3.1cm] (elf)
  {ELF 暂存区\\文件字节保持原样\\最多 512 KiB};
\node[os compute, text width=3.2cm, right=18mm of elf] (pass1)
  {第一遍：建立事实\\所有 PT\_LOAD 均合法\\入口落在可执行段};
\node[os memory, text width=3.1cm, right=18mm of pass1] (segments)
  {第二遍：建立映像\\复制文件区\\清零 BSS};
\node[os state, text width=3.1cm, below=14mm of pass1] (bootinfo)
  {BootInfo v2\\地址、长度、页表根\\栈与内存图};
\node[os control, text width=3.1cm, right=18mm of bootinfo] (entry)
  {RDI = BootInfo*\\16 字节栈约束\\调用内核入口};
\draw[os strong bus] (elf) -- (pass1);
\draw[os strong bus] (pass1) -- node[os label, above] {全体通过} (segments);
\draw[os strong bus] (segments) -- (entry);
\draw[os strong bus] (bootinfo) -- (entry);
\end{pdfdiagram}
\diagramnote{ELF 校验、映像写入和 ABI 交接是三个边界。第一遍让失败仍可回滚，
第二遍建立目标内存，BootInfo 再把启动事实交给内核。}

\section*{v0.5：内核必须替换临时描述符状态}

Stage 1 的 GDT 和页表只负责把加载器带进长模式，不应成为内核长期 ABI。
内核验证 BootInfo、BSS 和传入 CR3 后，构造自己的 GDT、104 字节 TSS 和
256 槽 IDT。TSS 保存 Ring 0 栈顶，并为双重故障、NMI 和机器检查指定三个
独立 IST 栈；IDT 的 0..31 项安装架构异常门。

异常入口分成硬件、汇编和 C++ 三层。CPU 保存 RIP、CS、RFLAGS 以及某些异常
的错误码；每向量汇编桩为“无硬件错误码”的异常补零并压入向量号；公共入口
再保存 15 个通用寄存器、清除 DF 并对齐 C++ 调用栈，形成 160 字节统一现场。
分发器允许受控 breakpoint 返回，其他异常进入无动态分配、不可重入的 panic。

\begin{pdfdiagram}
\node[os control, text width=3.3cm] (cpu)
  {CPU 硬件压栈\\RIP、CS、RFLAGS\\可选错误码};
\node[os compute, text width=3.3cm, right=18mm of cpu] (stub)
  {每向量汇编桩\\补齐错误码、压向量\\进入公共入口};
\node[os memory, text width=3.3cm, right=18mm of stub] (frame)
  {统一 160 B 现场\\RAX..R15、vector\\error、返回现场};
\node[os state, text width=3.3cm, below=14mm of stub] (dispatch)
  {C++ 分发\\INT3 可恢复\\其他异常记录并停机};
\draw[os strong bus] (cpu) -- (stub);
\draw[os strong bus] (stub) -- (frame);
\draw[os strong bus] (frame) |- (dispatch);
\end{pdfdiagram}
\diagramnote{统一异常帧由硬件保存部分、向量桩规范化部分和公共入口保存部分共同
组成。统一形状不表示所有异常具有相同恢复策略。}

这一里程碑是后续系统调用与用户异常隔离的硬件入口基础，却还不是用户边界。
当前 IDT 中 breakpoint 和 overflow 可以由 DPL3 软件触发，普通系统调用 ABI、
用户栈和用户地址验证随后在 v0.8 建立。正文的受控入口章节会在这个已实现基础上
继续推导，而不会把未来设计混入 v0.5 的完成声明。

\section*{v0.6：从启动内存图到内核拥有的页表}

Stage 1 通过 \code{fw\_cfg} 设备读取 \code{etc/e820}，把原始 20 字节条目转换
为项目的 24 字节物理内存图，并按物理基址排序。内核拒绝空表、零长度、溢出、
乱序、重叠以及受管范围内没有可用 RAM 的输入。这里的验证不是防止“数据不好看”，
而是阻止分配器把 MMIO、固件区或越界地址当成普通页帧。

页帧分配器当前管理低 64 MiB，共 16384 个 4 KiB 帧。每帧使用 2 bit 区分
unavailable、free、allocated 与 reserved。四种状态让“永远不可分配”“启动期
长期占用”和“动态所有者持有”具有不同失败语义；一位位图无法区分重复释放和
释放保留页。

\begin{longtable}{|L{0.19\linewidth}|L{0.25\linewidth}|L{0.27\linewidth}|L{0.19\linewidth}|}
\toprule
帧状态 & 谁可以建立 & 允许的后继 & 非法操作 \\
\midrule
unavailable & 内存图初始化 & 保持不变 & 分配、释放或动态保留 \\\tablerowrule
free & 可用 RAM 初始化或合法释放 & allocated、reserved & 再次释放 \\\tablerowrule
allocated & 分配器成功返回 & free；在完整预检后可转保留 & 重复分配、作为不可用区 \\\tablerowrule
reserved & 平台、内核映像、启动栈等保留过程 & 按明确生命周期解除 & 作为普通对象释放 \\
\bottomrule
\end{longtable}

内核不长期沿用 Stage 1 的 2 MiB 大页表，而是从页帧分配器取得新的四级
4 KiB 页表。零页和四个栈底 guard page 保持 not-present；内核 text 为 RX，
rodata 为 R/NX，data、BSS、栈与高半区 64 KiB 早期堆为 RW/NX。设置
\code{IA32_EFER.NXE} 和 \code{CR0.WP} 后，内核切换到新 CR3，再回读控制状态
并逐项查询映射。

\begin{pdfdiagram}
\node[os memory, text width=2.65cm] (map)
  {物理内存图\\范围、类型、顺序\\来自硬件配置设备};
\node[os compute, text width=2.65cm, right=12mm of map] (frames)
  {页帧分配器\\2 bit 所有权状态\\保留启动关键区};
\node[os memory, text width=2.65cm, right=12mm of frames] (tables)
  {新四级页表\\4 KiB 叶映射\\W\^X 与 guard};
\node[os control, text width=2.65cm, right=12mm of tables] (cr3)
  {启用 NX/WP\\写 CR3\\回读并查询映射};
\draw[os strong bus] (map) -- (frames);
\draw[os strong bus] (frames) -- node[os label, above] {分配并清零页表页} (tables);
\draw[os strong bus] (tables) -- node[os label, above] {全部结构完成后提交} (cr3);
\end{pdfdiagram}
\diagramnote{物理事实、页帧所有权、页表策略与处理器当前 CR3 是四层不同状态。
切换 CR3 是提交点，不替代切换前的完整构造与切换后的权限核验。}

写保护故障镜像让 Ring 0 写入一页 present、read-only 的高半区映射。预期
页故障错误码为 \code{0x3}，CR2 为精确测试地址。这个失败同时证明页存在、
访问类型为写、访问者为 supervisor，并证明 \code{CR0.WP} 没有被遗漏。相比
“查询页表项看起来只读”，真实故障把页表数据和处理器执行行为闭合起来。

\section*{v0.7：把外部事件接到内核状态机}

案例平台同时包含 8259A、I/O APIC 和本地 APIC。自研 ROM 没有传统 BIOS
替系统建立虚拟线模式，因此内核不能假设 8259A 的 INTR 已经能到达 CPU。
设备启动先清除并回读 \code{IA32\_APIC\_BASE[11]}，再初始化两片 8259A，
把 IRQ0..15 重映射到 IDT 32..47。

\begin{pdfdiagram}
\node[os memory, text width=2.5cm] (pit) {8254 PIT\\IRQ0\\递增 tick};
\node[os memory, text width=2.5cm, below=11mm of pit] (kbd)
  {i8042 + 键盘\\IRQ1\\扫描码状态机};
\node[os compute, text width=2.7cm, right=14mm of pit, yshift=-5.5mm] (pic)
  {8259A 主/从片\\重映射、mask\\ISR 与 EOI};
\node[os control, text width=2.7cm, right=14mm of pic] (idt)
  {IDT 32..47\\每向量 IRQ 桩\\公共硬件入口};
\node[os state, text width=2.7cm, right=14mm of idt] (runtime)
  {InterruptRuntime\\最小设备处理\\确认后返回};
\draw[os strong bus] (pit) -- (pic);
\draw[os strong bus] (kbd) -- (pic);
\draw[os strong bus] (pic) -- node[os label, above] {INTR} (idt);
\draw[os strong bus] (idt) -- (runtime);
\draw[os feedback] (runtime.south) |- node[os label, below] {按主/从来源发送 EOI} (pic.south);
\end{pdfdiagram}
\diagramnote{传统中断路径同时依赖设备、PIC 路由、IDT 入口和确认协议。IRQ 与
同步异常可以共用寄存器帧形状，但分发策略和返回前责任不同。}

第一次实现只重映射 PIC、开放 IRQ0 并执行 \code{STI}。当时 PIT 的 IRR bit0
已经置位，PIC mask 已开放，CPU 的 IF 也为 1，但向量 32 始终没有进入。增加
轮询时间或重写 PIT 都不能修复问题，因为本地 APIC 仍全局启用，传统 INTR
没有完成跨控制器路由。显式关闭并回读 LAPIC 全局使能后，IRQ0 与 IRQ1 才能
到达。这个真实故障说明：局部寄存器分别正确，不等于器件之间的连接契约已经
成立。

IRQ0 热路径只递增 64 位 tick，不格式化日志；IRQ1 读取一个扫描码并推进集合 1
解码器，语义事件在中断返回后的循环中消费。ATA 暂时保持同步单扇区 PIO，不开放
IRQ14，也没有 DMA 或通用块层。现有设备能力因此是后文长期请求、完成对象和
块层设计的起点，而不是这些高级机制的完成证明。

\section*{v0.8：用户代码第一次不再拥有整台机器}

v0.7 结束时，任何进入内核的 C++ 代码都以 Ring 0 权限执行。把应用程序直接
链接成另一个内核函数虽然容易，却不能阻止它改写页表、访问设备端口或因一次
非法指令拖垮整个系统。v0.8 因此先建立最小但真实的用户边界，而没有提前伪造
进程、线程或调度器。

用户程序由独立的 freestanding C++20 与 NASM 目标生成 AMD64
\code{ET_EXEC} ELF64，链接基址为 \code{0x40000000}。内核把 ELF 当作
不可信输入：先检查文件头、程序头、范围、对齐、W\^X、段间重叠、总页数和
可执行入口，再分配物理页并建立 U/S 映射。任何一步失败都逆序解除本轮映射和
页帧；验证失败不会留下半份用户布局。

\begin{longtable}{|L{0.25\linewidth}|L{0.22\linewidth}|L{0.43\linewidth}|}
\toprule
用户虚拟范围 & 权限 & 设计目的 \\
\midrule
\code{0..0xFFFF} & 不允许 & 捕获空指针附近和意外低地址访问 \\\tablerowrule
\code{0x40000000...} & 按 ELF 为 RX、R/NX 或 RW/NX & 装入首批用户代码、只读数据与可写数据 \\\tablerowrule
\code{0x7FFFFFFEA000} & not-present & 捕获用户栈向下越界 \\\tablerowrule
\code{0x7FFFFFFEB000..0x7FFFFFFEFFFF} & user RW/NX & 四个用户栈数据页 \\\tablerowrule
\code{0x00007FFFFFFF0000} & 仅作为栈顶边界 & 初始 RSP 指向半开区间末端，不建立叶映射 \\
\bottomrule
\end{longtable}

特权转换同时依赖描述符、页表和入口代码。GDT 增加选择子
\code{0x1B}/\code{0x23} 的 Ring 3 数据段与代码段；TSS.RSP0 指向一份带
guard 的独立 16 KiB 转换栈。内核保存原调用链后构造
\code{SS:RSP:RFLAGS:CS:RIP}，由 \code{iretq} 进入 CPL3。用户执行
\code{INT 0x80} 时，CPU 根据 TSS 换到可信内核栈并压入 176 字节特权现场；
普通调用恢复现场返回用户态，退出或用户异常则丢弃转换栈内容，回到原内核
调用者。

\begin{pdfdiagram}
\node[os memory, text width=3.0cm] (elf)
  {内嵌用户 ELF\\完整验证后\\建立 U/S 映像};
\node[os control, text width=3.0cm, right=15mm of elf] (enter)
  {保存内核调用链\\压五项特权帧\\\code{iretq} 到 CPL3};
\node[os state, text width=3.0cm, right=15mm of enter] (user)
  {用户 C++\\RX 代码 + RW/NX 栈\\只能使用受控 ABI};
\node[os compute, text width=3.0cm, below=14mm of enter] (syscall)
  {\code{INT 0x80}\\TSS.RSP0 换栈\\校验帧与用户指针};
\node[os state, text width=3.0cm, right=15mm of syscall] (result)
  {普通返回用户态\\或收束退出/异常\\恢复原内核调用链};
\draw[os strong bus] (elf) -- (enter);
\draw[os strong bus] (enter) -- (user);
\draw[os strong bus] (user) |- node[os label, right] {调用或故障} (syscall);
\draw[os strong bus] (syscall) -- (result);
\draw[os feedback] (result.north) -- node[os label, right] {普通调用 \code{iretq}}
  (user.south);
\end{pdfdiagram}
\diagramnote{用户边界不是一条降权指令。ELF 事务、U/S 页权限、TSS 换栈、
入口帧校验和退出/异常恢复共同决定用户错误是否能被限制在用户执行单元内。}

当前 ABI 使用 RAX 保存编号或结果，以 RDI、RSI 传递两个参数，
\code{WriteLog=1} 最多复制 160 字节，\code{ExitProcess=2} 终止当前用户
执行。内核先验证完整半开地址区间，再逐页检查 present 与 U/S，最后复制到
固定内核缓冲；非法指针、未知编号、过长写入和设备失败具有不同负错误值。正常
程序用六次系统调用验证错误路径、输出 \code{HELLO_FROM_RING3} 并以零退出。

完成证据还必须说明“用户出错时内核是否仍然活着”。专用镜像中的
\code{UD2} 产生用户向量 6；访问 \code{0x30000000} 产生向量 14、错误码
\code{0x4} 和同值 CR2。两条路径都只结束用户执行，内核恢复原调用链并继续
到 \code{READY}；同类 Ring 0 故障仍进入 panic。截断为 63 字节的用户 ELF
则在设备初始化和进入 Ring 3 之前被拒绝。

这个里程碑完成时仍只有一个同步用户执行单元，所有用户映像共用当前内核页表根，
结束后尚未回收用户页，也没有 PID、runqueue 或等待状态。随后完成的 v0.9
正是从这些限制出发，把已验证的保护边界装进具有独立生命周期的进程对象，再用
真实 PIT 引入定时器抢占和上下文切换。

\section*{v0.9 至 v1.0：四次失败怎样收束成一台系统}

v0.9 先处理“一个用户程序长期运行时，第二个程序何时获得处理器”。四槽 PCB
保存 PID、调度状态、tick 与派发统计；目标运行时另存页表根、176 字节用户现场
和每进程 Ring 0 栈。PIT 每四个 tick 重新选择候选者，真正切换时同时写 CR3、
TSS.RSP0 和恢复帧。三个 worker 使用相同 ELF 与相同虚拟 BSS 地址，只有独立
物理页才能让三份计数各自从零到达相同终值。

v0.10 随后处理“进程暂时不能推进时，为什么仍然占据 Ready 资格”。调度状态机
加入 Blocked 和具名等待原因；64 字节管道保存环形下标、占用量和两端关闭状态。
Try 只检查或提交，Wait 只改变调度资格，唤醒后用户包装必须重新 Try。生产者与
消费者传输 256 字节确定性载荷，31 字节消费块迫使部分传输与回绕；关闭写端让
消费者在读尽剩余数据后观察 EOF。

v0.11 再处理“内存中的成功怎样穿过重启”。2 MiB 磁盘先冻结启动链与文件系统的
所有权边界；文件系统区域用 superblock、两张位图、32 个 inode、64 字节目录项
和 1013 个数据块解释持久字节。修改先把 superblock 持久化为 Dirty，再写回
数据和元数据，最后恢复 Clean。该协议只能检测并拒绝半事务，不能执行日志重放；
同盘双启动与损坏后第三次停止分别验证正常持久化和失败边界。

v1.0 最后处理“各机制已经存在，普通用户程序怎样通过同一入口组合它们”。八槽
描述符表让 fd 0/1/2 引用标准输入、输出和错误，动态槽引用文件、目录或管道端点。
键盘字符从 i8042 经 IRQ1、Set 1 解码和 256 字节 FIFO 到达 Shell；Shell 阻塞
且没有其他 Ready 时，内核切回永久地址空间执行相邻的
\code{sti; hlt; cli}，由外部中断恢复。Shell 的十条命令只调用公开用户 ABI，
没有由内核或宿主直接解释命令文本。

\begin{pdfdiagram}
\node[os control, text width=2.8cm] (v09)
  {v0.9 进程\\独立 CR3/RSP0\\PIT 抢占与回收};
\node[os state, text width=2.8cm, right=11mm of v09] (v010)
  {v0.10 等待\\Blocked + 原因\\管道与关闭唤醒};
\node[os memory, text width=2.8cm, right=11mm of v010] (v011)
  {v0.11 持久化\\inode/目录/缓存\\Dirty/Clean};
\node[os compute, text width=2.8cm, right=11mm of v011] (v100)
  {v1.0 用户环境\\统一描述符\\键盘、空闲与 Shell};
\draw[os strong bus] (v09) -- node[os label, above] {不能推进} (v010);
\draw[os strong bus] (v010) -- node[os label, above] {需要跨重启} (v011);
\draw[os strong bus] (v011) -- node[os label, above] {需要用户组合} (v100);
\end{pdfdiagram}
\diagramnote{四个里程碑不是并列功能表。每一步都继承上一阶段的机器状态，再解决
一个尚不能表达的失败：执行资格、条件等待、持久身份和端到端用户工作流。}

v1.0 仍是明确有界的教学基线：单核、四进程、固定描述符容量、轮询 ATA、直接块
文件、最小 ASCII 终端和全部内建命令。它没有 fork/exec/wait、信号、作业控制、
通用等待队列、SMP、DMA、日志恢复、删除或 rename。正文会把已完成的最小结构
扩展到现代系统规模，同时继续把“项目真实拥有的能力”和“原理上需要继续增加的
机制”分开说明。

\section*{第二演进周期：从 v1.0 基线走向单核 Unix-like 系统}

\begin{statebox}
\textbf{状态边界。} 下列 v1.1 至 v2.0 是项目已经写入路线图的\textbf{规划阶段}，
不是当前仓库已经实现的功能。v1.0 的 73 项测试是本书可引用的现有完成证据；
后续版本只有在代码、失败路径和自动化验证同时落地后，才会进入上面的“已实现主线”。
\end{statebox}

第一演进周期解决了“系统能否从固件一直走到交互式用户环境”。第二周期不以继续
堆叠孤立命令为目标，而是要让资源、进程、文件和终端获得可以组合的 Unix 语义。
路线图把这一跨度拆成十个因果相邻的阶段；每一阶段都先补足下一阶段将依赖的状态
对象和生命周期，而不是提前暴露无法兑现的接口。

\begin{longtable}{|L{0.13\linewidth}|L{0.35\linewidth}|L{0.42\linewidth}|}
\hline
\rowcolor{BookBlueTint}
\textbf{规划阶段} & \textbf{要解除的 v1.0 限制} & \textbf{预期建立的契约} \\\tablerowrule
\endhead
v1.1 & 低 64 MiB、64 KiB 单调堆、四份静态内核栈和固定容量限制了对象生命周期。 &
管理全部 E820 可用页，引入可释放内核分配、动态 guard 栈、至少 64 个 PCB
和 64 槽描述符，并为引用计数与失败回滚奠基。 \\\tablerowrule
v1.2 & 当前直接块文件接口不能承载稳定名称空间、大文件或多种打开对象语义。 &
引入 vnode、mount、cwd、open-file description 与文件系统 v2，使路径解析、
对象操作和具体磁盘格式分层。 \\\tablerowrule
v1.3 & 用户映像仍由内核固定装入，进程之间没有父子关系、Zombie 与等待结果。 &
只由内核启动 PID1，再从 VFS 装载 ELF，建立 spawn/exec/wait、argv/envp
和退出结果的可追踪生命周期。 \\\tablerowrule
v1.4 & 创建新执行单元仍需完整重建地址空间，不能表达 Unix 的复制后分化。 &
用 fork、物理页引用计数和写时复制区分共享映射、私有修改及写故障时的
所有权转移。 \\\tablerowrule
v1.5 & Shell 还是内建命令集合，描述符不能完成继承、重定向和动态管道组合。 &
补齐 pipe、dup、close-on-exec、流水线与重定向，让外部 Shell 通过公开接口
编排独立的 \code{/bin} 程序。 \\\hline
\end{longtable}

\begin{longtable}{|L{0.13\linewidth}|L{0.35\linewidth}|L{0.42\linewidth}|}
\hline
\rowcolor{BookBlueTint}
\textbf{规划阶段} & \textbf{前一阶段留下的系统问题} & \textbf{预期建立的契约} \\\tablerowrule
\endhead
v1.6 & 外部程序已经可以组合，但没有时间等待、信号递送和可控终端会话。 &
建立 deadline、最小信号集、signal frame、进程组和 canonical 终端，使
前后台作业能够被等待、停止、继续或中断。 \\\tablerowrule
v1.7 & 地址空间仍以预先装入为主，缺页不能按映射来源恢复运行。 &
以 VMA 建立 mmap/brk、按需 ELF、受控栈增长和自研用户堆，使文件映射、
匿名页与用户运行时共享缺页协议。 \\\tablerowrule
v1.8 & ATA 路径仍以同步轮询为中心，Dirty 状态只能被检测和拒绝。 &
用 IRQ14、块请求队列、统一缓存和带 commit record 的日志明确完成、落盘、
重放、丢弃与故障传播边界。 \\\tablerowrule
v1.9 & 接口长期演化后，需要冻结兼容规则并系统清理越界、竞态和资源泄漏。 &
以 SYSCALL/SYSRET 冻结 ABI v2，加入 \code{/dev}、只读 \code{/proc}、审计、
泄漏快照与固定种子变异。 \\\tablerowrule
v2.0 & 各子系统分别完成仍不等于整机语义闭合。 &
以单核 Unix-like 工作流完成跨层集成、压力与恢复验证；SMP、网络和图形仍不纳入
本周期完成条件。 \\\hline
\end{longtable}

这条路线也给正文提供了清楚的阅读坐标：资源表对应对象生命周期，VFS 对应名称
与操作分派，进程树、fork 和 COW 对应执行身份与地址空间所有权，信号与终端对应
异步控制，按需分页对应可恢复的内存缺口，异步块层与恢复协议对应持久化完成。
因此，正文中的高级机制不是脱离项目的百科补充，而是对后续每个规划阶段所需契约
的提前推导；在项目真正实现之前，书中始终用“规划”而不是“已经支持”来表述。

\section*{怎样用案例阅读后面的深入章节}

重排后的正文首先跟随复位、模式转换、ELF 交接、异常、内存和设备里程碑，
再把已经完成的 v0.9 至 v1.0 分别放入调度、同步、文件系统和统一入口章节。
章节中的 Linux、现代文件系统或高级设备队列用于说明同一契约扩展到真实规模时
会增加什么状态；这些深入机制不会被反向写成项目 v1.0 已经实现的能力。

阅读每个项目节点时，至少保留四类证据：

\begin{enumerate}
\tightlist
\item \textbf{构造证据}：目标字节、描述符、页表项或对象字段已经写入预期位置。
\item \textbf{采用证据}：处理器或设备已经实际使用该状态，例如回读 CR3、触发
  页故障或接收真实 IRQ。
\item \textbf{失败证据}：损坏输入、超时或非法访问会进入预期错误终态，并禁止
  后续成功标记。
\item \textbf{跨层证据}：宿主观察、串口到达时间和来宾内部计数分别表达什么，
  不能用其中一种替代另一种。
\end{enumerate}

项目的最终目标不是让日志出现越来越多的 \code{READY}，而是让每个新增能力
都有清楚的输入、状态对象、所有权、提交边界、失败语义和验证方法。后面的深度
正文将始终沿这套问题展开。

\begin{pdfdiagram}
\node[os memory, text width=3.5cm] (boot)
  {已实现 v0.1--v0.4\\复位、装载、长模式\\ELF64 与 BootInfo};
\node[os control, text width=3.5cm, right=10mm of boot] (kernel)
  {已实现 v0.5--v0.8\\异常、内存与设备\\用户 ELF 和受控入口};
\node[os state, text width=3.5cm, right=10mm of kernel] (system)
  {已实现 v0.9--v1.0\\进程、同步与 IPC\\文件与用户环境};
\node[os compute, text width=3.5cm, below=9mm of boot] (bootdeep)
  {正文深入\\地址空间、程序映像\\控制权与交接事务};
\node[os compute, text width=3.5cm, below=9mm of kernel] (kerneldeep)
  {正文深入\\内核分配、请求队列\\中断、超时与恢复};
\node[os compute, text width=3.5cm, below=9mm of system] (systemdeep)
  {正文深入\\调度、并发、VFS\\网络、安全与持久化};
\draw[os strong bus] (boot) -- node[os label, right] {扩大规模与故障集合} (bootdeep);
\draw[os strong bus] (kernel) -- node[os label, right] {扩大并发与生命周期} (kerneldeep);
\draw[os strong bus] (system) -- node[os label, right] {扩展到现代系统} (systemdeep);
\end{pdfdiagram}
\diagramnote{项目里程碑提供全书的纵向控制权主线；正文不另起一套抽象体系，
而是在对应节点上增加规模、并发、延迟、隔离和故障恢复条件。}
